\section{Limitations}

The all-size claims concern frozen algebraic grammars and structural support objectives.  Structural support is not gate count, circuit depth, latency, qubit count, fault-tolerant overhead, error rate or quantum advantage.  Other TARE ranks, objectives outside the proved support-two cone, routing constraints and hardware cost vectors require new proofs.  The exact envelope in Eq.~\eqref{eq:envelope} is therefore not a production resource estimator.

The implementation evidence is extensive but dependent.  Thousands of rows are generated within a small number of formal families and do not constitute thousands of independent compiler domains.  The 6.34-billion-state calculation closes a finite local proof obligation whose arbitrary-size consequence comes from the analytic composition.  Computational scale improves checking; it does not by itself establish novelty, practical prevalence or external validity.

The state-preparation conclusions are finite.  Complete graphs are used only through $n=3$, and transfer is tested on one frozen 120-state $n=4$ panel.  The four-feature separator is exact on the smaller domain, but its mechanism attribution fails and its transfer result is adverse.  Other representations, gate sets, costs, folds or measurement-assisted synthesis may behave differently.

The remaining proof-route failures are not missing optional polish.  The single-pinner chain and its global $B''$ completeness implication remain unproved, while the two-coordinate chain premise is refuted.  Eq.~\eqref{eq:envelope} relies on a different per-block-permutation route.  All checks were performed within the same research programme and are not external proof certification or independent replication.
